% Folio: Drehen, im Bild (EP1), auf Deutsch. Begleitet Folio 01. Build: python engine/kl.py build <this file>
\documentclass[deutsch]{kaelearn}
\subject{Experimentalphysik I}
\folio{02}
\title{Drehen,\\im Bild}
\shorttitle{Drehen, im Bild}
\subtitle{Warum der Schwerpunkt so aussieht, wie sich Drehen anfühlt,\\und zehn Tricks, die Aufgaben kurz machen.}
\sources{Folio 01 · deine EP1-Mitschrift}
\readtime{2 Min. Überblick · 30 Min. Folio}
\filed{2026-11-09}
\plate{plate.png}{Das Orrery, zweite Nacht}
\kael{curious}

\newcommand\dd{\mathrm{d}}
\newcommand\ext{_{\mathrm{ext}}}
\newcommand\trick[2]{\par\medskip\needspace{5\baselineskip}\noindent{\tek\fontsize{11}{11}\selectfont\color{klhi}#1}\hspace{8pt}%
  {\fell\fontsize{13.5}{16}\selectfont\color{klhi}#2}\par\nopagebreak\smallskip\noindent\ignorespaces}
\tikzset{vec/.style={-{Stealth[length=2.2mm]},line width=.8pt},lab/.style={fill=klpage,inner sep=1.2pt}}

\begin{document}
\makecover

% ======================================================================
\section{Formeln, die man sehen kann}

\opening{D}{u kennst} die Formeln aus Folio 01. Was fehlt, ist das Gefühl, dass sie gar nicht anders aussehen \emph{könnten}. Eine Formel,
die du nur auswendig weißt, vergisst du in der Klausur in dem Moment, in dem die Aufgabe anders aussieht als die Übung. Eine Formel, die du
siehst, rekonstruierst du in zehn Sekunden.

Dieses Folio holt das nach: drei Bilder für den Schwerpunkt, eines für das Drehen, und zehn Tricks, die aus langen Aufgaben kurze machen.
Das erste Bild ist eine Wippe.

\begin{twominute}
\begin{itemize}
  \item \textbf{Schwerpunkt = Balancierpunkt.} Auf der Wippe gilt $m_1d_1 = m_2d_2$. Der Punkt, an dem sich die gewichteten Abstände
        aufheben, ist $\vec R_s = \sum m_i\vec r_i/M$: ein gewichteter Mittelwert, wie dein ECTS-Notenschnitt.
  \item \textbf{$\sum m_i\vec r_{si} = 0$} ist derselbe Satz: Vom Mittelwert aus heben sich die gewichteten Abweichungen auf. Darum
        balanciert ein Körper, im Schwerpunkt gestützt, in jeder Lage.
  \item \textbf{Der Punkt, der Newton gehorcht.} Nur wenn du die Orte mit den Massen gewichtest, fallen die inneren Kräfte heraus:
        $M\ddot{\vec R}_s = \vec F\ext$. Darum fliegt der Schwerpunkt eines geworfenen Hammers eine saubere Parabel.
  \item \textbf{Drehen ist Schieben in Verkleidung:} $\varphi,\omega,I,D,L$ statt $x,v,m,F,p$, solange die Achse fest ist.
  \item \textbf{$I$ misst, wo die Masse ist}, nicht wie viel: Jedes $\dd m$ muss mit $\omega r$ mitlaufen.
  \item \textbf{Drei echte Unterschiede:} $I$ hängt von der Achse ab; ein Drehmoment quer zu $\vec L$ dreht $\vec L$, statt es zu
        verlängern; $\vec L$ muss nicht parallel zu $\vec\omega$ sein.
  \item \textbf{Die stärksten Tricks:} Bezugspunkt auf die Unbekannten; Rollen ist Drehen um den Auflagepunkt; effektive Masse
        $m + I/r^2$; ein Loch ist negative Masse.
\end{itemize}
\end{twominute}

\begin{guessfirst}
\begin{enumerate}
  \item Warum ist $\sum_i m_i\vec r_{si} = 0$? Versuch es in einem Satz, ohne Formel.
  \item Eine lange Stange: Um welchen Faktor ist es schwerer, sie um ein Ende in Drehung zu versetzen als um ihre Mitte?
  \item Atwood-Maschine mit massiver Rolle (Vollzylinder, Masse $M$). Wie viel Masse „spürt“ die Beschleunigung von der Rolle: $M$, $M/2$ oder $0$?
\end{enumerate}
\end{guessfirst}

\begin{terms}
\term{Hebelgesetz}{law of the lever}{$m_1d_1 = m_2d_2$: die Wippe balanciert}
\term{gewichteter Mittelwert}{weighted mean}{$\sum w_ix_i\big/\sum w_i$}
\term{Schwerpunkt\-system}{centre-of-mass frame}{das System, in dem $S$ ruht; Gesamtimpuls dort null}
\term{Trägheits\-radius}{radius of gyration}{$k$ mit $I = Mk^2$}
\term{effektive Masse}{effective mass}{$m + I/r^2$: alles, was mitbeschleunigt werden muss}
\term{Momentanpol}{instantaneous centre}{der Punkt, der gerade ruht; alles dreht sich um ihn}
\end{terms}

% ======================================================================
\section{Warum der Schwerpunkt so aussieht}

\subsection{Bild 1: die Wippe}
\opening{L}{eg zwei} Massen auf ein leichtes Brett und such den Punkt, an dem du es mit einem Finger tragen kannst. Es balanciert, wenn sich
die Drehmomente der beiden Gewichte aufheben: $m_1g\,d_1 = m_2g\,d_2$. Das ist das Hebelgesetz, und das $g$ kürzt sich. Miss die Orte als
Zahlen entlang des Bretts, $x_1 < x_s < x_2$, also $d_1 = x_s - x_1$ und $d_2 = x_2 - x_s$:
\[
  m_1(x_s - x_1) = m_2(x_2 - x_s) \quad\Longleftrightarrow\quad x_s = \frac{m_1x_1 + m_2x_2}{m_1 + m_2}.
\]
Das ist die Schwerpunktformel. Niemand hat sie definiert, weil sie hübsch aussieht; sie ist die Antwort auf die Frage, wo der Finger
hin muss. Mit beliebig vielen Massen, links und rechts verteilt, lautet die Bedingung $\sum_i m_i(x_i - x_s) = 0$ (Massen links haben
negativen Abstand und drehen andersherum), also $x_s = \sum_i m_ix_i/M$. In drei Dimensionen gilt dasselbe für jede Koordinate.
\aside{Archimedes hat das Hebelgesetz um 250 v.\,Chr. bewiesen und damit Schwerpunkte von Dreiecken und Parabelsegmenten gefunden. Bewiesen, nicht geraten. Ich mochte ihn.}

Das ist ein \emph{gewichteter Mittelwert}, und den kennst du schon: Dein Notenschnitt wird mit den ECTS gewichtet. EP1 (8 ECTS) mit Note 1
und LA (4 ECTS) mit Note 4 geben $(8\cdot1 + 4\cdot4)/12 = 2$. Leg die Noten als Gewichte auf eine Zahlengerade, und die Wippe balanciert
genau beim Schnitt.

\begin{center}
\begin{tikzpicture}[x=2.3cm,y=1cm]
  \draw[klink,line width=1.3pt] (.55,0)--(5.45,0);
  \foreach \n in {1,...,5} \draw[kldim] (\n,0)--(\n,.12);
  \foreach \n in {2,3,5} \node[kldim,above] at (\n,.12) {\tek\fontsize{8}{8}\selectfont\n};
  \fill[pattern={dither50},pattern color=kldim] (2,-.03)--(1.82,-.72)--(2.18,-.72)--cycle;
  \draw[kldim] (2,-.03)--(1.82,-.72)--(2.18,-.72)--cycle;
  \fill[pattern={dither25},pattern color=klink] (.78,.04) rectangle (1.22,1.64); \draw[klhi] (.78,.04) rectangle (1.22,1.64);
  \fill[pattern={dither25},pattern color=klink] (3.78,.04) rectangle (4.22,.84); \draw[klhi] (3.78,.04) rectangle (4.22,.84);
  \node[above,align=center] at (1,1.64) {\fell\itshape EP1, Note 1\\[-2pt]\fell\itshape 8 ECTS};
  \node[above,align=center] at (4,.84) {\fell\itshape LA, Note 4\\[-2pt]\fell\itshape 4 ECTS};
  \draw[kldim,<->,line width=.5pt] (1,-.95)--(1.97,-.95); \node[kldim,below] at (1.5,-.95){$1$};
  \draw[kldim,<->,line width=.5pt] (2.03,-.95)--(4,-.95); \node[kldim,below] at (3,-.95){$2$};
  \node[klhi,below] at (2,-1.35){\fell\itshape Schnitt $2$};
\end{tikzpicture}
\end{center}
\plotcaption{Links $8\cdot1$, rechts $4\cdot2$: Die Wippe balanciert beim Notenschnitt. Ein Schwerpunkt ist ein Notenschnitt, bei dem
Massen die Rolle der ECTS spielen.}

\subsection{Warum \texorpdfstring{$\sum m_i\vec r_{si} = 0$}{die Summe null ist}}
Das ist die Frage, an der alles hängt, darum vier Antworten, jede aus einer anderen Richtung.

\emph{Die Rechnung.} $\sum m_i\vec r_{si} = \sum m_i(\vec r_i - \vec R_s) = \sum m_i\vec r_i - M\vec R_s = M\vec R_s - M\vec R_s = 0$. Eine Zeile,
und sie zeigt: Der Satz ist keine Entdeckung über den Körper. Er ist die Definition von $\vec R_s$, umgestellt.

\emph{Der Mittelwert.} Die gewichteten Abweichungen vom gewichteten Mittelwert summieren sich immer zu null. Im Notenbeispiel:
$8\cdot(1-2) + 4\cdot(4-2) = -8 + 8 = 0$. Das liegt nicht an deinen Noten; der Mittelwert ist \emph{die} Zahl, um die herum das gilt.
Läge er woanders, wären die Abweichungen auf einer Seite zu schwer.

\emph{Die Wippe.} Um $S$ hat die Schwerkraft das Drehmoment $\sum_i\vec r_{si}\times m_i\vec g = \big(\sum_i m_i\vec r_{si}\big)\times\vec g$.
Null, egal wohin $\vec g$ zeigt, also in jeder Lage des Körpers. Genau das heißt „balanciert“.

\emph{Die Bewegung.} Leite ab: $\sum_i m_i\dot{\vec r}_{si} = 0$. Relativ zum Schwerpunkt hat der Körper keinen Gesamtimpuls. Was sich um $S$
dreht, dreht sich so, dass sich die Impulse in alle Richtungen aufheben. Darum ist der Gesamtimpuls $\vec p = M\vec v_s$, und darum zerfällt
die Energie sauber in $\frac12Mv_s^2$ plus Drehung: Der Kreuzterm ist $\vec v_s\cdot\sum m_i\dot{\vec r}_{si} = 0$.

\why{Warum balanciert ein Körper, den du genau im Schwerpunkt stützt, in jeder Lage, und warum kann das kein anderer Punkt?}{Um einen
Stützpunkt $P$ ist das Schweremoment $\big(\sum m_i\vec r_{Pi}\big)\times\vec g$. Für $P = S$ ist die Klammer null, also das Moment für
jede Lage. Für $P\ne S$ ist die Klammer $M(\vec R_s-\vec r_P)\ne0$, und sobald $\vec g$ nicht parallel dazu steht, kippt der Körper. Ein
Hammer balanciert auf dem Finger nur an einer Stelle des Stiels, nahe am Kopf.}

\why{Warum liegt der Schwerpunkt eines Rings dort, wo gar keine Masse ist?}{Ein Mittelwert muss keiner der gemittelten Werte sein: Einen
Schnitt von 2,5 hat niemand als Note bekommen. Beim Ring hat jedes Stück ein gegenüberliegendes, die beiden heben sich in
$\sum m_i\vec r_{si}$ auf, und übrig bleibt die leere Mitte.}

\subsection{Bild 2: der Punkt, der Newton gehorcht}
Wirf einen Hammer. Er dreht sich wild, jeder seiner Punkte fliegt eine verschlungene Kurve. Bis auf einen. Schreib Newton für jedes Stück
und addiere:
\[
  \sum_i m_i\ddot{\vec r}_i = \sum_i\vec F_i^{\,\mathrm{ext}} + \sum_{i\ne j}\vec F_{ij} = \vec F\ext ,
\]
denn innere Kräfte kommen in Paaren $\vec F_{ij} = -\vec F_{ji}$. Links steht die zweite Ableitung von $\sum m_i\vec r_i = M\vec R_s$. Also
$M\ddot{\vec R}_s = \vec F\ext$: Der Schwerpunkt bewegt sich, als säße die ganze Masse dort und alle äußeren Kräfte griffen dort an.
Beim Wurf ist $\vec F\ext = M\vec g$, und er fliegt eine Parabel, ganz gleich wie der Hammer sich dreht.

Und jetzt das Entscheidende: \emph{Diese Gewichtung ist erzwungen.} Nimm irgendeinen gewichteten Mittelwert $\sum c_i\vec r_i$ mit
$\sum c_i = 1$. Seine Beschleunigung ist $\sum_i(c_i/m_i)\big(\vec F_i^{\,\mathrm{ext}}+\sum_j\vec F_{ij}\big)$. Damit sich jedes Paar
innerer Kräfte aufhebt, egal wie es aussieht, muss $c_i/m_i = c_j/m_j$ für alle Paare gelten, also $c_i = m_i/M$. Die Schwerpunktformel ist
der einzige Mittelwert der Orte, der von den inneren Kräften nichts merkt.

\begin{center}
\begin{tikzpicture}[scale=.92]
  \draw[kllow,line width=.6pt] plot[domain=0:1,samples=320,variable=\s] ({-5.5+11*\s+1.2*cos(1080*\s+20)},{-1.6+12*\s*(1-\s)+1.2*sin(1080*\s+20)});
  \draw[kldim,line width=.6pt] plot[domain=0:1,samples=200,variable=\s] ({-5.5+11*\s-.4*cos(1080*\s+20)},{-1.6+12*\s*(1-\s)-.4*sin(1080*\s+20)});
  \draw[klhi,line width=1.2pt] plot[domain=0:1,samples=80,variable=\s] ({-5.5+11*\s},{-1.6+12*\s*(1-\s)});
  \foreach \s in {0,.25,.5,.75,1}{
    \pgfmathsetmacro\cx{-5.5+11*\s}\pgfmathsetmacro\cy{-1.6+12*\s*(1-\s)}\pgfmathsetmacro\a{1080*\s+20}
    \draw[klink,line width=1.1pt] ({\cx+1.2*cos(\a)},{\cy+1.2*sin(\a)})--({\cx-.4*cos(\a)},{\cy-.4*sin(\a)});
    \fill[klink] ({\cx+1.2*cos(\a)},{\cy+1.2*sin(\a)}) circle (2.2pt);
    \fill[klhi] ({\cx-.4*cos(\a)},{\cy-.4*sin(\a)}) circle (3.8pt);
  }
  \node[klhi,lab] at (0,2.05) {\fell\itshape Schwerpunkt: Parabel};
  \node[kldim,lab,anchor=west] at (-6.6,-3.0) {\fell\itshape leichtes Ende $m$: Schleifen};
  \node[kldim,lab,anchor=east] at (6.6,-3.0) {\fell\itshape schweres Ende $3m$: Wellen};
\end{tikzpicture}
\end{center}
\plotcaption{Eine geworfene Hantel ($m$ und $3m$, der Schwerpunkt ein Viertel vom schweren Ende). Fünf Momentaufnahmen. Die Enden schlingern,
der Schwerpunkt zieht eine glatte Parabel, weil sich die Stabkräfte in der Summe aufheben.}

\anim{schwerpunkt}{Zwei Würfe derselben Hantel. Beim ersten schlingern die Enden, beim zweiten auch die Stabmitte; nur der Punkt ein
Viertel vom schweren Ende fliegt glatt. Dann balanciert die Hantel genau dort, und Newton zeigt, warum die Stabkraft herausfällt.}

\begin{keystone}{Der Schwerpunkt}{centre of mass}
Genau ein Punkt erfüllt $\sum_i m_i(\vec r_i - \vec R_s) = 0$, nämlich $\vec R_s = \frac1M\sum_im_i\vec r_i$. Dieser Punkt
\begin{enumerate}
  \item balanciert im homogenen Schwerefeld in jeder Lage,
  \item bewegt sich nach $M\ddot{\vec R}_s = \vec F\ext$, egal was der Körper sonst tut,
  \item teilt Energie und Drehimpuls sauber: $E = \frac12Mv_s^2 + E_{\text{um }S}$ und $\vec L = M\vec R_s\times\vec v_s + \vec L_s$.
\end{enumerate}
\smallskip\noindent{\color{kldim}\pix{Beweis in einem Atemzug}}\quad (1) $\big(\sum m_i\vec r_{si}\big)\times\vec g = 0$. (2) Innere Kräfte
heben sich paarweise auf. (3) Jeder Kreuzterm enthält $\sum m_i\vec r_{si}$ oder $\sum m_i\dot{\vec r}_{si}$, und beide sind null.
\end{keystone}

% ======================================================================
\section{Wie sich Drehen anfühlt}

\opening{S}{olange die} Achse fest ist, ist Drehen nichts Neues: Es ist Schieben, mit anderen Buchstaben. Jede Formel für einen Massenpunkt
hat einen Zwilling, und du darfst übersetzen, ohne nachzudenken.

\begin{panel}[unbreakable]{Tabelle}{Übersetzen}{translation and rotation}
\renewcommand\arraystretch{1.3}
\begin{tabular}{@{}p{24mm}p{30mm}p{58mm}@{}}
  {\color{kldim}\pix{Schieben}} & {\color{kldim}\pix{Drehen}} & {\color{kldim}\pix{Was du spürst}}\\
  $x$, $v$, $a$ & $\varphi$, $\omega$, $\dot\omega$ & Winkel statt Weg; außen läuft jeder Punkt mit $v = \omega r_\perp$\\
  $m$ & $I = \int r_\perp^2\,\dd m$ & wie schwer es ist, in Drehung zu kommen\\
  $F$ & $D = r_\perp F$ & Türklinke außen: wenig Kraft, viel Wirkung\\
  $p = mv$ & $L = I\omega$ & wie schwer es ist, die Drehung zu stoppen\\
  $F = \dot p$ & $D = \dot L$ & Drehmoment ändert den Drehimpuls\\
  $\frac12mv^2$ & $\frac12I\omega^2$ & Energie, die in der Drehung steckt
\end{tabular}
\end{panel}

\subsection{Trägheitsmoment: wo, nicht wie viel}
Dreht sich ein Körper mit $\omega$, muss jedes $\dd m$ mitlaufen, und zwar außen schneller: $v = \omega r_\perp$. Masse weit draußen zählt
doppelt. In der Energie: Sie muss schneller laufen, $v = \omega r_\perp$, und die Energie geht mit $v^2$. Im Drehimpuls: Sie trägt mehr
Impuls, $\dd m\,\omega r_\perp$, und das an einem längeren Hebel, $r_\perp$. Beide Wege enden bei $r_\perp^2$. Eine Stange um ihr Ende statt um ihre Mitte: $\frac13M\ell^2$ statt $\frac1{12}M\ell^2$, viermal so träge, obwohl die Masse
dieselbe ist. Gleiche Masse, andere Achse, anderes $I$.

Ein Bild, das hilft, ist der \emph{Trägheitsradius} $k$: $I = Mk^2$. Er ist der Radius eines dünnen Rings derselben Masse mit demselben $I$.
Vollzylinder $k = R/\sqrt2\approx0{,}71\,R$, Vollkugel $k = \sqrt{2/5}\,R\approx0{,}63\,R$, Stange um die Mitte $k = \ell/\sqrt{12}\approx0{,}29\,\ell$.
Eine Vollkugel benimmt sich beim Drehen, als säße ihre ganze Masse auf einem Ring mit 63\,\% ihres Radius.
\aside{Die Erde hat $I\approx0{,}33\,MR^2$, nicht $0{,}4$: Innen ist sie dichter. So wiegt man einen Eisenkern, den niemand je gesehen hat.}

\subsection{Drei Stellen, an denen die Übersetzung bricht}
\emph{Erstens: $I$ gehört zur Achse, nicht zum Körper.} Eine Masse hat ein $m$; ein Körper hat für jede Achse ein anderes $I$. Darum fragst
du bei jeder Aufgabe zuerst: um welche Achse?

\emph{Zweitens: $\vec L$ ist ein Pfeil, und ein Drehmoment quer dazu dreht ihn.} $\dd\vec L = \vec D\,\dd t$ zeigt in Richtung des
Drehmoments. Steht das senkrecht auf $\vec L$, wird $\vec L$ nicht länger, nur schwenkt es. Der Kreisel kippt nicht, er wandert seitwärts.
Merkregel: Die Spitze von $\vec L$ läuft dem Drehmoment nach.

\needspace{16\baselineskip}
\begin{center}
\begin{tikzpicture}[scale=1.25]
  \draw[kldim,dashed,line width=.5pt] (4,0) arc (0:24:4);
  \draw[vec,klhi,line width=1.1pt] (0,0)--(4,0) node[below]{$\vec L$};
  \draw[vec,klink] (4,0)--(4,1.2) node[right]{$\dd\vec L = \vec D\,\dd t$};
  \draw[vec,klhi] (0,0)--(4,1.2) node[above left]{$\vec L + \dd\vec L$};
  \draw[kldim] (1.3,0) arc (0:16.7:1.3); \node[kldim] at (1.72,.2){$\dd\varphi$};
\end{tikzpicture}
\end{center}
\plotcaption{Ein kleiner Zuwachs senkrecht zu $\vec L$ ändert die Länge nur in zweiter Ordnung, $|\vec L+\dd\vec L| = \sqrt{L^2+\dd L^2}\approx L$,
aber die Richtung in erster: $\dd\varphi = \dd L/L$. Viele solche Schritte ergeben eine Kreisbewegung der Spitze, die Präzession $\Omega = D/L$.
(Die Zeichnung übertreibt $\dd\vec L$.)}

\emph{Drittens: $\vec L$ muss nicht parallel zu $\vec\omega$ sein.} Bei symmetrischen Körpern, die sich um eine Symmetrieachse drehen, ist es
so; sonst nicht (Folio 01, Teil V: Hauptachsen, Unwucht). Prüfe es, statt es anzunehmen.

\why{Warum fühlt sich ein schnell drehendes Rad „steif“ an, wenn du seine Achse kippen willst?}{Kippen heißt: die Richtung von $\vec L$
ändern. Bei großem $L$ ist schon eine kleine Richtungsänderung ein großes $\dd\vec L$, und das verlangt ein großes Drehmoment
$\vec D = \dot{\vec L}$. Dazu weicht die Achse seitwärts aus, denn $\dd\vec L$ zeigt in Richtung deines Drehmoments, nicht in Richtung deiner
Hand. Probier es mit einem Fahrradrad: Es ist das seltsamste Gefühl, das die Mechanik zu bieten hat.}

% ======================================================================
\section{Zehn Tricks}

\opening{J}{eder dieser} Tricks spart in irgendeiner Aufgabe eine halbe Seite. Lern sie nicht als Liste; lern bei jedem, \emph{wann} er
greift. Das steht jeweils im ersten Satz.

\trick{01}{Bezugspunkt auf die Unbekannten}
Wenn eine Kraft unbekannt ist und dich nicht interessiert (Gelenk, Lager, Auflagepunkt), stell die Momentengleichung um genau ihren
Angriffspunkt auf. Ihr Hebelarm ist null, sie verschwindet. Leiter an der Wand: um den Fuß, und von drei Unbekannten bleibt eine. Erlaubt ist
das um $S$ oder einen raumfesten Punkt, in der Statik um jeden Punkt.

\trick{02}{Rollen ist Drehen um den Auflagepunkt}
Wenn ein Körper rollt (oder irgendeine ebene Bewegung einen Punkt hat, der gerade ruht), ist die Bewegung in diesem Moment eine reine
Drehung um diesen Punkt, den Momentanpol. Jeder Punkt bewegt sich senkrecht zu seiner Linie dorthin, mit $\omega$ mal Abstand. Oben am Rad:
$2v_s$, ohne Rechnung. Energie in einem Zug: $E = \frac12I_P\omega^2$ mit $I_P = I_s + MR^2$. Bei einer rutschenden Leiter liegt der Momentanpol
dort, wo sich die Senkrechten auf die Geschwindigkeiten der beiden Enden schneiden.

\trick{03}{Effektive Masse}
Wenn ein System nur eine Koordinate $x$ hat und alles über Seile ohne Schlupf oder Rollen ohne Gleiten an $\dot x$ hängt, dann ist die
kinetische Energie $\frac12m_{\mathrm{eff}}\dot x^2$ mit
\[
  m_{\mathrm{eff}} = \sum m + \sum\frac{I}{r^2},
\]
wobei $r$ der Radius ist, an dem $v = \omega r$ gilt. Leistet nur die antreibende Kraft $F$ Arbeit, folgt aus
$\frac{\dd}{\dd t}\big(\frac12m_{\mathrm{eff}}\dot x^2\big) = F\dot x$ sofort $a = F/m_{\mathrm{eff}}$. Ein Drehkörper fühlt sich an wie die
Zusatzmasse $I/r^2$, an der die Schwerkraft nicht zieht.

\begin{panel}[unbreakable]{Tabelle}{Effektive Masse}{a = Antrieb / effektive Masse}
\renewcommand\arraystretch{1.35}
\begin{tabular}{@{}p{56mm}p{26mm}l@{}}
  {\color{kldim}\pix{Aufbau}} & {\color{kldim}\pix{Antrieb}} & {\color{kldim}\pix{effektive Masse}}\\
  Vollzylinder rollt die schiefe Ebene hinab & $Mg\sin\alpha$ & $M + \frac{\frac12MR^2}{R^2} = \frac32M$\\
  Atwood, Rolle als Vollzylinder $M$ & $(m_2-m_1)\,g$ & $m_1+m_2+\frac12M$\\
  Jo-Jo, Achsradius $r$ & $mg$ & $m + I_s/r^2$\\
  Klotz $m_1$ auf glattem Tisch, gezogen von hängender Masse $m_2$ über eine Rolle $I$, $r$ & $m_2\,g$ & $m_1+m_2+I/r^2$
\end{tabular}
\end{panel}

Erste Zeile: $a = \frac23g\sin\alpha$. Zweite: $a = (m_2-m_1)g/(m_1+m_2+\frac12M)$, dieselbe Antwort wie in Folio 01, ohne eine einzige
Seilkraft. Die Rolle „wiegt“ für die Beschleunigung $M/2$; das war Frage 3.

\why{Warum fügt ein Drehkörper genau $I/r^2$ zur Masse hinzu?}{Seine Energie ist $\frac12I\omega^2$, und mit $\omega = v/r$ wird daraus
$\frac12\,\frac{I}{r^2}\,v^2$: dieselbe Form wie $\frac12mv^2$ mit $m = I/r^2$. Für die Beschleunigung zählt nur, wie viel Energie pro
Geschwindigkeit$^2$ im System steckt. Ob sie in Translation oder Drehung steckt, ist der Bewegungsgleichung egal.}

\trick{04}{Ein Loch ist negative Masse}
Wenn ein Körper eine Aussparung hat, rechne den vollen Körper und zieh das Loch als Körper mit negativer Masse ab. Das geht für den
Schwerpunkt (Summen sind linear) und für $I$ (Integrale sind additiv; das Loch mit Steiner um dieselbe Achse). Scheibe, Radius $R$, Masse $M$
ohne Loch, mit kreisförmigem Loch vom Radius $R/2$ bei $x = R/2$: Das Loch hätte die Masse $M/4$, also
\[
  x_s = \frac{M\cdot0 - \frac M4\cdot\frac R2}{M - \frac M4} = -\frac R6 .
\]
Der Schwerpunkt rückt vom Loch weg. Auch der dicke Hohlzylinder ist so in einer Zeile: $\frac12M_2R_2^2 - \frac12M_1R_1^2 = \frac12M(R_1^2+R_2^2)$.

\begin{center}
\begin{tikzpicture}[scale=1.05]
  \fill[pattern={dither25},pattern color=klink,even odd rule] (0,0) circle (1.6) (.8,0) circle (.8);
  \draw[klink,line width=.7pt] (0,0) circle (1.6); \draw[kldim,line width=.6pt] (.8,0) circle (.8);
  \draw[kldim,dashed,line width=.4pt] (-2,0)--(2.2,0);
  \fill[klink] (0,0) circle (1.4pt) node[above=2pt,lab]{$O$};
  \fill[klhi] (-.267,0) circle (2.2pt) node[below=3pt,lab]{$S$};
  \node[kldim] at (.8,.38){\fell\itshape Loch};
  \node[anchor=west,kldim] at (2.5,.3){\begin{minipage}{45mm}\fell\itshape Volle Scheibe minus Loch.\\ Der Schwerpunkt liegt bei $-R/6$.\end{minipage}};
\end{tikzpicture}
\end{center}

\trick{05}{Die Symmetrie-Summe}
Wenn ein Körper in alle drei Richtungen gleich aussieht (Kugel, Hohlkugel, Würfel), addiere die drei Achsen:
$I_x+I_y+I_z = 2\int r^2\,\dd m$, und jedes ist ein Drittel davon. Vollkugel: $\int r^2\,\dd m = \frac35MR^2$, also $\frac25MR^2$. Hohlkugel:
$\int r^2\,\dd m = MR^2$, also $\frac23MR^2$. Keine Kugelkoordinaten, kein $\sin\theta$.

\trick{06}{Flache Körper: senkrechte Achsen}
Wenn die ganze Masse in einer Ebene liegt, ist $I_z = I_x + I_y$ für die Achse senkrecht zur Ebene. Scheibe um einen Durchmesser:
$\frac12\cdot\frac12MR^2 = \frac14MR^2$. Quadratische Platte, Kante $a$: $I_z = \frac16Ma^2$, also $\frac1{12}Ma^2$ um jede Achse in der Ebene
durch die Mitte, auch die Diagonale, denn die vierzählige Symmetrie macht alle diese Achsen gleich.

\trick{07}{Steiner nur über den Schwerpunkt}
Wenn die Achse nicht durch $S$ geht: $I_A = I_s + Md^2$, wobei eine der beiden Achsen durch $S$ gehen \emph{muss}. Von $A$ nach $B$ immer über
$S$. Und weil $Md^2\ge0$, ist die Achse durch $S$ die leichteste aller parallelen Achsen. Stange: Mitte $\frac1{12}$, Ende
$\frac1{12}+\frac14 = \frac13$, der Faktor 4 aus Frage 2.

\trick{08}{Die $k$-Probe}
Schreib jedes Ergebnis als $I = kMR^2$ mit $R$ = größter Abstand von der Achse. Dann muss $0\le k\le1$ sein, denn $r_\perp\le R$ überall; $k = 1$
nur für den Ring. Kommt $k>1$ heraus, hast du einen Fehler. Und jede Formel mit $I$ muss für $I\to0$ das Ergebnis für einen Massenpunkt geben
(Rolle ohne Masse, Klotz statt Kugel).

\trick{09}{Drehimpuls um die Wirkungslinie der Unbekannten}
Wenn eine Kraft kurz und heftig oder unbekannt ist (Stoß, Gleitreibung bis zum Rollen), nimm den Drehimpuls um einen raumfesten Punkt
auf ihrer Wirkungslinie. Dort hat sie kein Moment, und $L$ bleibt erhalten, obwohl Impuls und Energie es nicht tun. Kugel im Stab: um das
Lager. Billardkugel vom Gleiten zum Rollen: um einen Punkt auf dem Tisch, $v = v_0/(1+k)$.

\trick{10}{Energie für „wie schnell“, Kräfte für „wie stark“}
Wenn nach einer Geschwindigkeit zwischen zwei Lagen gefragt ist und nichts Energie verheizt (Rollen ohne Gleiten, reibungsfreie Lager),
nimm den Energiesatz, mit $\frac12(M+I_s/R^2)v^2$. Wenn nach einer Kraft gefragt ist (Seilkraft, Reibung, Lagerkraft), brauchst du die
Bewegungsgleichungen, $M\vec a_s = \vec F$ und $D = I\dot\omega$. Gleitreibung oder ein Stoß im Spiel: Energie ist verloren, nimm Trick 09.

% ======================================================================
\begin{triggers}
\trigger{Seil, Rolle oder Rollen, nur eine Koordinate}{effektive Masse: $a = F_{\mathrm{Antrieb}}\big/\big(\sum m + \sum I/r^2\big)$}
\trigger{Körper mit Loch oder Aussparung}{Loch als negative Masse, für $\vec R_s$ und für $I$}
\trigger{Ebene Bewegung, Geschwindigkeit eines Punktes gesucht}{Momentanpol: $v = \omega\cdot$Abstand, senkrecht zur Verbindung}
\trigger{Kugel, Hohlkugel, Würfel: $I$ gesucht}{$I_x+I_y+I_z = 2\int r^2\,\dd m$}
\trigger{Dünne Platte}{$I_z = I_x + I_y$}
\trigger{Ein Ergebnis mit $I$}{$k$-Probe: $0\le k\le1$; $I\to0$ gibt den Massenpunkt}
\trigger{Schwerpunkt mehrerer Teile}{gewichteter Mittelwert der Teil-Schwerpunkte}
\end{triggers}

\needspace{.42\textheight}
\marginkael{delighted}
\begin{deck}
\card{Warum}{Warum ist $\sum_i m_i\vec r_{si} = 0$?}{$S$ ist der gewichtete Mittelwert; die gewichteten Abweichungen davon heben sich auf. Physikalisch: Dort balanciert der Körper.}
\card{Warum}{Warum fliegt gerade der Schwerpunkt eine Parabel?}{Nur bei Gewichten $\propto m_i$ fallen die inneren Kräfte aus der Summe: $M\ddot{\vec R}_s = M\vec g$.}
\card{Trigger}{Seil, Rolle oder Rollen, eine Koordinate. Wozu greifst du?}{$a = F_{\mathrm{Antrieb}}/(\sum m + \sum I/r^2)$.}
\card{Trigger}{Scheibe mit Loch, Schwerpunkt gesucht. Wozu greifst du?}{Volle Scheibe plus Loch mit negativer Masse.}
\card{Def}{Trägheitsradius: $I = $ \{\{$Mk^2$\}\}; Vollkugel $k = $ \{\{$\sqrt{2/5}\,R\approx0{,}63\,R$\}\}.}{Der Radius des Rings mit gleicher Masse und gleichem $I$.}
\end{deck}

\end{document}
